% STATUS: verified
\chapter{The model layer}\label{ch:models}

Every other chapter of this part follows one path on which a model is asked something. This
chapter describes what all of those paths share: how a model is reached at all, what one call
looks like on the wire, how the shape of its answer is enforced, what is recorded about it,
what it costs and who pays. It goes from the outside in. \Cref{sec:models:door} states the
principle of one door to the models; \cref{sec:models:providers} the providers and the ids of
models; \cref{sec:models:generate,sec:models:embed} the two routes of the assistant service,
one for generated text and one for embeddings; \cref{sec:models:client} the client the API
calls them with, and its error handling. The second half is about money and limits: what is
recorded for a call (\cref{sec:models:usage}), the budgets (\cref{sec:models:budgets}), the
provider's prompt cache and what it did to the design (\cref{sec:models:cache}), rate limits
(\cref{sec:models:ratelimits}), and the model each step runs on by default, with the measured
reason for the choice (\cref{sec:models:defaults}).


\section{One door to the models}\label{sec:models:door}

The system runs as several services (\cref{ch:arch}). Two of them matter here. The
\term{API} is the Django service that holds all data and, through its \term{worker}, runs the
background \term{jobs} --- reading an act, indexing a document, making a report. The
\term{assistant service} (\code{legalops-assistant}: Node~24, Express~5 and version~7 of
the AI~SDK, as \file{package.json} in \file{assistant/app/} declares them) is the only process
that ever opens a connection to a model. The API does not call a model: when a job step
needs one, it sends an HTTP request to the assistant service, which makes the call and
returns the answer. The chat assistant lives inside the assistant service and reaches the
models through the same code (\cref{ch:chat}).

The division is deliberate, and \file{docs/markdown/en/architecture.md} states it as two design rules:
\enquote{Django never calls a model provider} and \enquote{the assistant owns no data}.
Three things follow from it.

\begin{itemize}
  \item \textbf{Providers and keys are configured once.} Which providers exist, what their
    addresses and keys are, and how each one is spoken to is written in one file of one
    service (\file{llm/provider.js}). Adding a provider or changing the key touches nothing
    else.
  \item \textbf{The API never holds a key.} The settings of the API name the assistant
    service (\code{ASSISTANT_URL}), the shared secret of the door between them and the id of
    the embedding model --- no provider's address and no provider's key. The process that
    parses uploaded files and serves the public routes cannot leak what it does not have.
  \item \textbf{Every call is logged the same way.} Whatever asked --- a chat turn, a job
    step, an embedding batch, a test run in the prompt editor --- the result is one row in the
    table \code{llm_calls} with the same columns (\cref{sec:models:usage}). Costs, budgets
    and the usage page read that one table.
\end{itemize}

The assistant service, in turn, stores nothing: prompts, conversations and usage rows live in
the API's database, and the assistant service reads and writes them through the API.

\paragraph{The internal door.} The routes the API calls are under \code{/internal/llm/}
(\file{routes/internal.js}; \cref{tab:models:routes}). They exist only on the isolated
network between the containers: the public reverse proxy forwards only paths under
\code{/api/} to the services, so a path that begins with \code{/internal/} never reaches them
from outside. Every request must also carry the header \code{X-Internal-Secret}.
The gate (\file{middleware/internal.js}, \code{requireInternal}) compares it with the
configured secret in constant time and runs before the body of the request is read; an empty
configured secret refuses every call with 403~\code{INTERNAL_ONLY} instead of letting
everything through. The API guards its own internal routes the same way
(\file{assistant/api/internal_views.py}, \code{require_internal}). A request body on this door
may be 16~MB; the public chat route, which takes one message, allows 512~kB.

\begin{table}[H]
\caption{The routes of the assistant service that the API calls}\label{tab:models:routes}
\begin{tabularx}{\textwidth}{@{}L{58mm}Y@{}}
\toprule
\hd{Route} & \hd{What it does} \\
\midrule
\route{POST}{/internal/llm/generate} & One generation: plain text, or an object that follows a JSON schema (\cref{sec:models:generate}). \\
\route{POST}{/internal/llm/embed} & The embeddings of a batch of texts (\cref{sec:models:embed}). \\
\route{GET}{/internal/llm/models} & The configured providers and the gateway's models with their prices; the source of the price list (\cref{sec:models:usage}) and of the model picker in the prompt editor. \\
\bottomrule
\end{tabularx}
\end{table}

\Cref{fig:models:path} shows one call from end to end. The rest of the chapter walks along
it.

\begin{figure}[H]
\centering
\begin{tikzpicture}[x=1mm, y=1mm]
  % the API
  \node[svc, text width=27mm] (step) at (14,0) {a job step\\{\scriptsize or a route a person waits at}};
  \node[svc, text width=27mm] (rp) at (14,-18) {\code{run_prompt}\\{\scriptsize finds the prompt version, fills its variables}};
  \node[svc, text width=27mm] (gen) at (14,-37) {\code{gateway.generate}\\{\scriptsize retries; waits out a rate limit}};
  \node[chk, text width=27mm] (rec) at (14,-56) {\code{record_call}\\{\scriptsize \code{cost_of}: each part of the input at its price}};
  \node[store] (calls) at (14,-74) {\code{llm_calls}};
  \node[store] (pv) at (49,-18) {prompt\\versions};
  \node[store] (prices) at (49,-56) {\code{model_prices}};
  % the assistant service
  \node[chk, text width=28mm] (gate) at (89,-37) {\code{requireInternal}\\{\scriptsize the shared secret}};
  \node[svc, text width=28mm] (route) at (89,-56) {the generate route\\{\scriptsize asks for a stream, reads it to its end}};
  \node[chk, text width=28mm] (check) at (89,-74) {the object is checked against the schema};
  % outside
  \node[ext, text width=17mm] (gw) at (130,-56) {the gateway};
  \node[ai, text width=17mm] (model) at (130,-74) {the model};

  \draw[flow] (step) -- (rp);
  \draw[flow] (rp) -- (gen);
  \draw[flow] (pv) -- node[lbl, above] {active} (rp);
  \draw[call] (gen) -- node[lbl, above] {POST \code{/internal/llm/generate}} (gate);
  \draw[flow] (gate) -- (route);
  \draw[call] (route) -- node[lbl, above] {Responses API} (gw);
  \draw[call] (gw) -- (model);
  \draw[flow] (route) -- (check);
  \draw[flow] (gen) -- node[lbl, right, pos=.45] {answer and five token counts} (rec);
  \draw[flow] (prices) -- node[lbl, above] {prices} (rec);
  \draw[flow] (rec) -- node[lbl, right] {one row} (calls);

  \begin{scope}[on background layer]
    \node[lane, fit=(step)(calls)(pv)(prices)] (laneapi) {};
    \node[lane, fit=(gate)(route)(check)] (laneas) {};
  \end{scope}
  \node[lanelabel, anchor=south west] at (laneapi.north west) {The API (Django): the worker or the web process};
  \node[lanelabel, anchor=south west] at (laneas.north west) {The assistant service};
\end{tikzpicture}
\actorlegend
\caption{The path of one model call. A job step calls \code{run_prompt}, which takes the
instructions and the model from the active prompt version; the API's client posts the request
through the internal door; the assistant service asks the gateway for a stream, reads it to
its end, checks the object and answers whole. Usage is written in one place: after the answer
has come back, \code{record_call} prices the call and inserts one row into
\code{llm_calls}. A dashed arrow is a request together with its answer.}\label{fig:models:path}
\end{figure}


\section{Providers and model ids}\label{sec:models:providers}

A \term{provider} is whoever runs models behind an API. A model is named everywhere in the
system by one string of the form \texttt{provider:model}, for example
\texttt{gateway:azure/gpt-6.1-sol}: a prefix that selects the provider, a colon, and the id the
provider itself gives the model. This string is what a prompt version stores in its
\code{model} field (\cref{ch:prompts}), what a usage row records and what the price list is
keyed by.

\subsection{The registry}

\file{llm/provider.js} holds the registry: one entry per prefix, saying whether the provider
has its credentials, how its client is built (once, on first use) and how it makes a language
model and an embedding model (\cref{tab:models:providers}). \code{parseModelId} splits a
string at its first colon, but only when what stands before the colon is a known prefix; any
other string is read as an id of the gateway, so \model{azure/gpt-6.1-sol} and
\texttt{gateway:azure/gpt-6.1-sol} are the same model. The answer of a call always names the
model in the full form (\code{qualifiedModel}), because the API prices the call by it.

\begin{table}[H]
\caption{The provider registry of the assistant service}\label{tab:models:providers}
\begin{tabularx}{\textwidth}{@{}L{21mm}L{37mm}YL{28mm}@{}}
\toprule
\hd{Prefix} & \hd{Configured when} & \hd{Language models} & \hd{Embeddings} \\
\midrule
\texttt{gateway:} & \code{AI_GATEWAY_KEY} is set (address: \code{AI_GATEWAY_URL}) & The gateway's Responses API, every call stateless & The gateway's \code{/embeddings} \\
\texttt{openai:} & \code{OPENAI_API_KEY} is set & OpenAI, through its own client of the SDK & yes \\
\texttt{anthropic:} & \code{ANTHROPIC_API_KEY} is set & Anthropic, through its own client of the SDK & none: asking for one is an error \\
\texttt{ollama:} & \code{OLLAMA_BASE_URL} is set & An Ollama server, as OpenAI-compatible chat models & yes \\
\bottomrule
\end{tabularx}
\end{table}

A provider without credentials fails loudly: \code{providerFor} throws, and the route answers
400~\code{UNKNOWN_MODEL} before any call is made. A call is never handed to another provider
in its place, and a string that names no model at all is refused in the same way rather than
sent on to fail in the middle of a stream.

Of the four, this installation uses one. \file{docs/markdown/en/architecture.md} records the decision ---
the faculty's gateway is the only model provider --- and the network enforces it: the only
model host on the allowlist of outgoing connections is the gateway's (\cref{ch:security}).
The other three prefixes are optional: each is switched on by setting its key or its
address, and the compose file of this installation passes none of them to the assistant
service.

\subsection{The gateway and its wire format}

The \term{gateway} is the faculty's Bifrost gateway (default address
\texttt{https://ai.knf.vu.lt/v1}). It speaks the wire formats of OpenAI and passes each call on
to a model deployed behind it; the ids of its models begin with \code{azure/}. The key is sent
in two headers at once, as \code{x-bf-vk} and as a bearer token in \code{Authorization}.

The registry builds two clients on that one key, because the two kinds of call use two
different formats:

\begin{itemize}
  \item \textbf{Language models are called on the Responses API} (\code{/v1/responses}), not
    on the older chat-completions format. The reason is in the header of
    \file{llm/provider.js}: the gateway's newer models (GPT-6 Astra, GPT-6.1 Sol) accept
    function tools only there, and the chat assistant always sends its tools.
  \item \textbf{Embeddings are called on} \code{/v1/embeddings}, through the SDK's
    OpenAI-compatible client.
\end{itemize}

The list of models and prices is read from \code{/v1/models}.

\paragraph{Stateless calls.} The Responses API can keep a conversation on the provider's
side and let the next request refer to it. The system never uses that.
\code{gatewayLanguageModel} wraps every gateway model so that each request carries the option
\code{store} set to \code{false}: neither the gateway nor the provider behind it keeps what
was asked and answered, and a request always contains everything the model is to read. For a
structured generation that is simply one request and one answer. For the chat it means that
each step of a turn sends the earlier steps back in full (\cref{ch:chat}).

\paragraph{Reasoning models.} Some models \enquote{reason} before they answer: they produce
internal tokens that are billed as output and are not part of the answer. The SDK recognises
such models by OpenAI's own ids, but the gateway's ids carry a deployment prefix
(\code{azure/}) that the SDK does not expect. \code{isReasoningModel} therefore applies the
same rule to the name after the last slash.

\begin{codebox}{How a reasoning model is recognised, and what changes}
A gateway model counts as a reasoning model when its name is \code{gpt-N} or
\code{gpt-N.M} with $N \ge 5$ and any suffix --- except a whole-number version whose suffix
begins with \code{chat} --- or when it belongs to the \code{o} series (the letter and a
number). For such a model the wrapper does three things:
\begin{itemize}
  \item it marks the model as reasoning (\code{forceReasoning}), so that the SDK sends the
    instructions as the \enquote{developer} message and asks for the model's reasoning in
    encrypted form --- which a later step of a stateless chat turn has to hand back;
  \item it removes \code{temperature} and \code{topP} from the request, because a reasoning
    model takes neither;
  \item it leaves \code{store} off, as for every model.
\end{itemize}
Every model the shipped prompts name (\cref{sec:models:defaults}) is a reasoning model by
this rule. The temperature stored on a prompt version is therefore not sent with the shipped
defaults; it takes effect only when an administrator moves a prompt to a model that is not a
reasoning model.
\end{codebox}

\paragraph{The fallback models.} Two settings of the assistant service name the model to use
when a request names none: \code{AI_CHAT_MODEL} (default \texttt{gateway:azure/gpt-5.5}) for
generation and \code{AI_EMBED_MODEL} (default \texttt{gateway:azure/text-embedding-3-large})
for embeddings. Every shipped prompt version names its own model, so the first is a safety
net. The second is the model of every search query's vector; the stored vectors are made with
the API's own setting of the same name (\code{AI_EMBED_MODEL} in \file{settings.py}), which
the callers send with every batch, so the two settings must name the same model
(\cref{sec:models:embed}).


\section{The generate route}\label{sec:models:generate}

\route{POST}{/internal/llm/generate} makes one generation. It is the route behind every
structured answer in the system: the obligations of a batch of provisions, the verdict on an
act, the findings for fifteen obligations, an action plan. This section goes through the
request, the way the model is read, the limits in time, the schema and the failures.

\subsection{What the API sends}

The request is a JSON object. \code{readGenerateBody} validates it field by field before any
model is involved (\cref{tab:models:request}); a field that fails is a 400 with a code of its
own, and nothing reaches the model.

\begin{table}[H]
\caption{The fields of a generate request}\label{tab:models:request}
\begin{tabularx}{\textwidth}{@{}L{27mm}L{45mm}Y@{}}
\toprule
\hd{Field} & \hd{Rule} & \hd{What it carries} \\
\midrule
\code{prompt} & required, a non-empty string (\code{INVALID_PROMPT}) & The data of this one call: the provisions to read, the obligations to judge. It is sent as the single user message. \\
\code{system} & optional string (\code{INVALID_SYSTEM}) & The instructions: the text of the prompt version with its variables filled in. \\
\code{shared} & optional string (\code{INVALID_SHARED}) & Data that every call of one step carries unchanged, for example the company's profile and documents in a report. Placed behind the instructions (\cref{sec:models:cache}). \\
\code{model} & optional string, trimmed & The model id of the prompt version; without one, \code{AI_CHAT_MODEL}. \\
\code{schema} & optional JSON object, not a list (\code{INVALID_SCHEMA}) & The JSON schema the answer must follow. Without it the answer is plain text. \\
\code{schemaName} & optional string, cut to 64 characters; default \code{result} & The name the schema is sent under. \\
\code{temperature} & optional number from 0 to 2 (\code{INVALID_TEMPERATURE}) & The sampling temperature of the prompt version; dropped for a reasoning model. \\
\code{maxOutputTokens} & optional whole number from 1 to 128\,000 (\code{INVALID_MAX_TOKENS}) & The output cap of the prompt version, when it sets one. \\
\bottomrule
\end{tabularx}
\end{table}

\begin{aibox}{What the model receives in a structured generation}
The route joins \code{system} and \code{shared} with a blank line and sends the result as
the \emph{instructions} of the request; \code{prompt} follows as one user message; the
schema, when there is one, goes as a strict response format under its name. That is all. There
is no conversation history, there are no tools, and nothing of an earlier call is referred to:
whatever the model is to know about the law, the company or the task is in these three texts.
What they contain on each path is the subject of the chapters that follow.
\end{aibox}

\subsection{Read as a stream, answered whole}

The route does not ask the provider for a finished answer. It asks for a stream
(\code{streamText} of the SDK), reads the stream to its end inside the assistant service, and
only then answers the API with the whole result in one JSON object:

\begin{codeblock}
{ "model": "gateway:azure/gpt-6.1-sol",
  "text": null,  "object": { ... },      one of the two is null
  "finishReason": "stop",
  "usage": { "inputTokens", "cachedInputTokens", "cacheWriteTokens",
             "outputTokens", "reasoningTokens" },
  "durationMs": 187000 }
\end{codeblock}

The API therefore sees an ordinary request that takes long. The stream exists for one reason,
stated in the header of \file{routes/internal.js}: a slow model writes a long answer for many
minutes, and Node drops a response whose first byte takes longer than five minutes to arrive.
A stream begins at once, so the wait for the first byte is short however long the model
writes.

\begin{lessonbox}{Measured, 10 November 2026}
On that day the GPT-6 models of the gateway wrote about 70 tokens a second. A dense batch of
provisions sent for the extraction of obligations took four to six minutes; one verified call
ran 338~seconds and wrote 24\,000 output tokens. A call of that length cannot be made as a
single response under a five-minute limit on the first byte, which is why the generation is
read as a stream and why the ceilings below are fifteen minutes.
\end{lessonbox}

A stream can also fail after it has begun --- the provider reports an error as an event inside
it. The SDK would end such a stream quietly, which would look like a shorter answer. The route
keeps the first error it sees (\code{onError}) and throws it once the stream is over:
\enquote{a stream that reports an error after it began is a failed call, never a shorter
answer}, in the words of the route's header. A test in \file{tests/internal.test.mjs} pins it
(\enquote{no partial answer reaches the worker}).

\subsection{Timeouts on both sides}

The two sides of the door have matching ceilings (\cref{tab:models:timeouts}). The assistant
service gives one generation fifteen minutes; the API waits half a minute longer, so that the
assistant's refusal arrives before the API gives up and the reason is not lost.

\begin{table}[H]
\caption{The time limits of a model call}\label{tab:models:timeouts}
\begin{tabularx}{\textwidth}{@{}L{42mm}L{50mm}Y@{}}
\toprule
\hd{Constant} & \hd{Where} & \hd{Value and meaning} \\
\midrule
\code{GENERATE_TIMEOUT_MS} & assistant service, \file{config.js} & 900\,000~ms: one generation. Past it the provider call is aborted and the route answers a 502 that names the timeout. \\
\code{GENERATE_TIMEOUT_S} & API, \file{assistant/gateway.py} & 930~s: how long \code{generate} waits for the assistant service. \\
\code{EMBED_TIMEOUT_MS} & assistant service, \file{config.js} & 60\,000~ms: one embedding batch, retries included. \\
\code{timeout} of \code{embed} & API, \file{assistant/gateway.py} & 120~s: twice the assistant's ceiling, so that a hung call becomes a 502 the client retries instead of running on after the API gave up. \\
\code{maxRetries} & assistant service, both routes & 2: the SDK repeats a failed provider call twice (a spent budget is asked once). \\
\bottomrule
\end{tabularx}
\end{table}

\subsection{The schema, and an answer that does not fit it}

A step that needs structured output sends a JSON schema. The schemas are owned by the code
that parses the answers --- each pipeline module defines its own --- and
\file{prompts/schemas.py} lists them by prompt (\code{SCHEMAS}: 21 prompts, each with the
module, the constant and the name the schema is sent under), so that a test run in the prompt
editor uses the very schema the pipeline does. The smallest of them shows the shape:

\begin{codefile}{\file{corpus/citations.py} --- \code{JUDGE_SCHEMA}, sent under the name \code{verdicts}}
JUDGE_SCHEMA = {
    "type": "object", "additionalProperties": False, "required": ["verdicts"],
    "properties": {"verdicts": {"type": "array", "items": {
        "type": "object", "additionalProperties": False, "required": ["ref", "verdict", "note"],
        "properties": {"ref": FILLED,
                       "verdict": {"type": "string", "enum": ["match", "mismatch", "unclear"]},
                       "note": FILLED}}}},
}
\end{codefile}

Every schema is written for the gateway's \term{strict mode}: every property is listed as
required, \code{additionalProperties} is \code{false}, and a value that may be absent is
declared as a nullable type instead of an optional key. The schema is sent as a strict response
format, which asks the provider to hold the model's output to it: no missing key, no key of
the model's own, no status outside the listed ones.

\paragraph{Required is not filled.}\label{par:models:filled} Strict mode makes every key
required, and that promises the key only: \code{{"type": "string"}} is satisfied by an empty
string and an array by an empty list. What an answer cannot do without is therefore marked in
the schema itself (the header of \file{prompts/schemas.py} states the rule):

\begin{itemize}
  \item a text that must carry content is \code{FILLED}, a string with \code{minLength}~1 ---
    the \code{ref} and the \code{note} above;
  \item a list with a known size says so with \code{minItems} and \code{maxItems}: the eight
    parts of a policy, the five levels of a criteria table, the one provision at least behind
    an obligation.
\end{itemize}

A field the prompt allows to be empty stays a plain string: a condition, a deadline, a
paragraph number. The line between the two is the prompt's own wording, because a mark the
prompt does not back is worse than none --- the decoder forces \emph{something} into a
filled field, and a model with nothing to say fills it with a dash. The gateway enforces
these keywords while the model writes and the assistant service's validator checks them
again; both were tried on all three models before the schemas were marked.

\begin{lessonbox}{A mark that changed the answers, 17 November 2026}
Some fields are needed in some answers only: a recommendation when a finding is open, the
questions when an act is left uncertain, a note when an obligation is doubted. A schema can
say that too, as two shapes of the item under \code{anyOf} --- a closed finding whose
recommendation may be empty, an open one whose recommendation is filled --- and the gateway
accepts it. It was built, and the measurement decided against it. With the two shapes the gap
analysis agreed with 64 and 65 of the evaluation's 71 labelled cases; with the schema as it
was, on the same day, with 67 and 68; with one shape and the plain marks of this section,
with 69 and 69 (\cref{ch:control}). Two cases moved in both runs of the two-shape form and
in none of the other four, each towards the open status: a finding the labels hold to be met
was answered as partial, another as unknown; what else differed changed from run to run. Nothing in the two shapes forbids the right answer; the form of
the schema is part of what the model reads, and it shifted a judgement. The three schemas
that had the form have one shape again, and what they could have enforced is asked by the
prompt alone. A mark is added to a schema only where the evaluation shows that the answers
are no worse with it.
\end{lessonbox}

\begin{codebox}{The answer is checked twice}
\code{schemaOutput} in \file{routes/internal.js} passes the schema to the provider as a strict
response format \emph{and} builds a validator from the same schema, which the parsed answer
must pass inside the assistant service before it is returned. The header gives the reason: a
provider without strict mode may drift, and the API reads the object by key, so a renamed key
would silently mean \enquote{no findings}. An answer that cannot be parsed as JSON or does not
match the schema fails the call; it never reaches the API as a result. A schema the validator
cannot read is sent unchecked, with a warning in the log, and then only the provider's strict
mode stands behind it.
\end{codebox}

A failed check carries no HTTP status of a provider, so the route answers it as
502~\code{GENERATION_FAILED}, which the API's client treats as transient and asks once more
(\cref{sec:models:client}). A schema the gateway itself refuses --- one that is not
strict-shaped --- comes back as the gateway's own 4xx and is not retried: asking again cannot
help.

\subsection{The path card}

Every structured generation in the system takes this path. The cards of the later chapters say
what is particular to each prompt; this one says what is common to all.

\begin{pathcard}{A structured generation}{run_prompt}
  \pcrow{Trigger}{A job step of the worker, or a route a person waits at, calls
    \code{run_prompt} in the API (\file{assistant/llm.py}) with a prompt's slug and the data. The test run of the
    prompt editor calls \code{gateway.generate} directly with the draft's text.}
  \pcrow{Input}{Three texts: the instructions (the active prompt version, its
    \code{{PLACEHOLDER}} variables filled); optionally a shared block that all calls of the
    step carry; the data of this call. Plus the JSON schema and its name, the model id, the
    temperature and the output cap of the prompt version. The body may be 16~MB; batch sizes
    and size caps belong to each path.}
  \pcrow{Prompt and model}{The prompt named by the slug; the model of its active version,
    unless the caller passes another. A call made for a job step is \emph{patient}: it waits
    out a rate limit. A reasoning model is sent no temperature.}
  \pcrow{Output}{With a schema, \code{object}: the parsed JSON. Without one, \code{text}.
    Always \code{model}, \code{finishReason}, \code{durationMs} and the five token counts of
    \code{usage}.}
  \pcrow{Checked by code}{The request, field by field. The schema is enforced by the
    provider's strict mode and checked again in the assistant service. An error inside the
    stream fails the call. Everything that concerns the \emph{content} of the answer is checked
    by the calling path.}
  \pcrow{Stored}{One row in \code{llm_calls}: kind \code{job_step} (\code{prompt_test} for a
    test run), the purpose, the model,
    the five token counts, the cost, the duration, the job step, the user charged and the
    prompt version. The answer itself is stored by the caller. The system keeps no cache of
    answers at this layer; the provider caches the instructions (\cref{sec:models:cache}).}
  \pcrow{On failure}{The SDK repeats a failed provider call twice, a spent budget never. The
    API's client repeats a transient failure once; a job step waits out a rate limit for up to
    fifteen minutes; a lost connection is retried five times. Then \code{GatewayError} reaches
    the caller (\cref{tab:models:errors}). A failed call leaves no usage row.}
  \pcrow{Cost}{Priced part by part from the gateway's price list and charged to the user the
    caller names, or to nobody (\cref{tab:models:charged}). Measured on 10--11 November 2026 on
    the sample company: a gap-analysis call of 59\,000 input tokens cost \$0.18 when it filled
    the provider's cache and about \$0.045 when it read from it.}
  \pcrow{Human control}{None at this layer. An administrator chooses the text and the model of
    every prompt (\cref{ch:prompts}) and the daily budget of every account; what a person
    reviews in the answer is said in the card of each path.}
  \pcrow{Measured by}{Tests, not evaluations: \file{tests/internal.test.mjs},
    \file{tests/provider.test.mjs} and \file{tests/usage.test.mjs} in the assistant service
    (the wire, the schema check, failures, usage); \file{tests/test_prompts.py} in the API
    (pricing, retries, patience, the price list). Nothing at this layer measures a prompt.}
\end{pathcard}


\section{The embed route}\label{sec:models:embed}

An \term{embedding} is a vector of numbers computed from a text, such that texts with similar
meaning get vectors that lie close together. The system stores one for every chunk of the
legal corpus, every obligation and every chunk of a company's documents, and computes one for
every search query; \cref{ch:corpus,ch:company} say what is done with them. All of them come
from \route{POST}{/internal/llm/embed}.

\paragraph{The request.} A list \code{texts} of 1 to 256 strings (\code{MAX_EMBED_TEXTS};
otherwise 400~\code{INVALID_TEXTS} or \code{TOO_MANY_TEXTS}) and optionally a \code{model};
without one the route uses \code{AI_EMBED_MODEL}. An empty string would make some providers
fail the whole batch, so it is replaced by a single space. The SDK's \code{embedMany} sends
the batch with at most four parallel calls and two retries, under the ceiling
\code{EMBED_TIMEOUT_MS}.

\paragraph{The model and its width.} The embedding model is
\model{text-embedding-3-large} of the gateway. Its vectors are asked for at
\textbf{1536 dimensions}, not at the model's full width, because that is the width of every
vector column in the database (\code{EMBED_DIMENSIONS} in \file{corpus/models.py} and
\file{companies/models.py}, repeated in \file{routes/internal.js}). The route passes the
\code{dimensions} option only to models of the \code{text-embedding-3} family, the ones that
accept it. As a second line, \code{fitDimensions} cuts any longer vector to its first 1536
values and renormalises it --- the same thing the option does --- so that a provider that
ignored the option still fits. A shorter vector is passed on and the database column refuses
it.

Because every stored vector must come from the same model to be comparable, changing
\code{AI_EMBED_MODEL} means embedding everything again (the job \code{corpus.reembed};
\cref{ch:corpus}).

\paragraph{The answer.} \code{embeddings} (one vector per text, in the order sent),
\code{model}, \code{dimensions}, \code{usage.tokens} and \code{durationMs}.

\paragraph{Batching.} The route's limit is 256 texts; the callers stay well below it
(\cref{tab:models:embedcallers}). Each batch is one call and one usage row.

\begin{table}[H]
\caption{Who asks for embeddings, and in what batches}\label{tab:models:embedcallers}
\begin{tabularx}{\textwidth}{@{}L{43mm}YL{40mm}@{}}
\toprule
\hd{Caller} & \hd{Text embedded} & \hd{Batch; purpose recorded} \\
\midrule
\file{corpus/ingest.py}, \code{embed_chunks} & A chunk of an act, opened by the act's short title, the provision's path and its heading & 64 (\code{EMBED_BATCH}); \code{corpus.embed} \\
\file{corpus/obligations.py}, \code{embed_obligations} & An obligation's title and text, and its English title and text where it has them & 32 obligations, at most 64 texts; \code{obligations.embed} \\
\file{corpus/obligations.py}, \code{embed_english} & The English title and text of obligations that lack that vector & 64; \code{obligations.embed_en} \\
\file{companies/indexing.py}, \code{index_document} & A window of a company document, opened by the document's title and the section heading & 64; \code{documents.embed} \\
\file{corpus/search.py}, \file{companies/search.py} & The query of a search & 1; not recorded \\
\bottomrule
\end{tabularx}
\end{table}

\begin{pathcard}{An embedding call}{gateway.embed}
  \pcrow{Trigger}{Ingesting or re-embedding an act (jobs \code{corpus.ingest},
    \code{corpus.reembed}); extracting obligations, or an administrator editing one; indexing
    an uploaded document (job \code{documents.index}); a report that meets a document not yet
    indexed or obligations without an English vector; every search of the corpus or of a
    company's documents.}
  \pcrow{Input}{A list of texts, at most 256 per call (\code{MAX_EMBED_TEXTS}); the callers
    send at most 64, or one query (\cref{tab:models:embedcallers}). Each text is the passage together
    with what places it: the act and provision, or the document and section.}
  \pcrow{Prompt and model}{No prompt. \texttt{gateway:azure/text-embedding-3-large}
    (\code{AI_EMBED_MODEL}), asked for 1536 dimensions.}
  \pcrow{Output}{\code{embeddings}, one vector of 1536 numbers per text in the order sent;
    \code{model}; \code{dimensions}; \code{usage.tokens}; \code{durationMs}.}
  \pcrow{Checked by code}{The texts are validated and empty ones replaced by a space. A longer
    vector is cut to 1536 values and renormalised. Corpus ingestion refuses an answer whose
    number of vectors differs from the number of texts. The database column refuses a vector
    of another width.}
  \pcrow{Stored}{The vector columns \code{legal_chunks.embedding},
    \code{obligations.embedding} and \code{obligations.embedding_en},
    \code{company_document_chunks.embedding}. A query vector is used once and not kept. One
    \code{llm_calls} row of kind \code{embed} per batch, except for a query.}
  \pcrow{On failure}{Two retries in the SDK; a batch that outlives 60~s is aborted and
    answered as a 502. The API's client retries a transient failure three times (after 2, 6
    and 18~s). A search whose query cannot be embedded runs on keywords alone. An indexing job
    is attempted three times.}
  \pcrow{Cost}{\$0.13 per million tokens (the gateway's price list, 10 November 2026).
    The embeddings of an upload are charged to the uploader; corpus work to nobody; a query is
    not recorded. Measured: ingesting 24 new acts on a copy of the database cost \$0.08 on 12 November 2026;
    re-embedding everything was put at some \$0.30 on 28 September 2026.}
  \pcrow{Human control}{None.}
  \pcrow{Measured by}{\file{tests/internal.test.mjs} and \file{tests/embed-timeout.test.mjs}
    (the route, the 1536 dimensions, the ceiling); the comparison of the two embedding models
    of 28 September 2026 (\cref{tab:models:embedcompare}).}
\end{pathcard}


\section{The API's client}\label{sec:models:client}

On the API's side one module talks to the assistant service, \file{assistant/gateway.py}, and
one function on top of it is what the pipelines call, \code{run_prompt} in
\file{assistant/llm.py}. The module is called \enquote{gateway} because it is the API's
gateway to the models; it should not be confused with the faculty's gateway, which it never
contacts.

\subsection{The two calls}

\code{gateway.generate} builds the request of \cref{tab:models:request} from its arguments,
leaves out those that are \code{None}, posts it with the shared secret and returns a
\code{GenerateResult}: the text or the object, the model, the finish reason, the five token
counts and the duration. \code{gateway.embed} does the same for a list of texts and returns an
\code{EmbedResult}. A third function, \code{list_models}, reads the model list and answers an
empty one when the assistant service is down, so that nothing that only wants prices can fail
on it.

Both calls go through \code{_post}, a loop that sends the request and decides what to do
with a failure (\cref{sec:models:errors}). A failure that is final is raised as
\code{GatewayError} with three things the caller can act on: the HTTP status (0 for a network
failure), a machine code, and the message of the assistant service, which carries the
provider's own words. The text of a network error names internal hosts and ports; it goes to
the log, and the error a user may read in a report or a job says what happened in words
(\enquote{The AI service did not answer within 930~s}).

\subsection{run\_prompt: a stored prompt over data}

\code{run_prompt} is, in the words of its header, \enquote{the one way a pipeline step talks
to a model}. Its signature shows what a caller decides and what it does not:

\begin{codeblock}
run_prompt(slug, data, *, shared=None, schema=None, schema_name="result",
           variables=None, job_step_id=None, user_id=None, purpose=None, model=None)
\end{codeblock}

It does four things in order.

\begin{enumerate}
  \item \textbf{It finds the prompt version.} \code{active_version(slug)}
    (\file{prompts/store.py}) returns the version an administrator has made active for that
    prompt. If there is none, it raises \code{PromptNotConfigured} and the step fails with
    that message: a step never runs on instructions that nobody can see.
  \item \textbf{It fills the variables.} \code{render} replaces each placeholder of the form
    \code{{NAME}} (capital letters, digits, underscores) that the caller supplied a value for;
    a report passes \code{LANGUAGE}, \code{TODAY} and \code{COMPANY_NAME}. A placeholder
    without a value stays in the text as written, so a mistyped name in the editor shows up in
    the instructions instead of vanishing.
  \item \textbf{It calls the model.} The rendered text is the instructions; \code{data} is the
    user message; \code{shared}, the schema and its name are passed through; the model, the
    temperature and the output cap come from the prompt version (a caller may override the
    model). The call is \emph{patient} exactly when \code{job_step_id} is given.
  \item \textbf{It records the call} with \code{record_call} (\cref{sec:models:usage}),
    including the id of the prompt version, \enquote{so a report can always be traced to the
    exact instructions that produced it}.
\end{enumerate}

It returns a \code{PromptRun}: the object or the text, the raw result and the version used.
The caller chooses the data, the schema, who is charged and under which job step the call is
filed; the text of the instructions and the model are never in the caller's code. That is what
lets an administrator change either without a deployment (\cref{ch:prompts}).

\subsection{Shared data, and calls in parallel}

Two helpers shape how a step issues many calls.

\code{parallel_map(fn, items, workers=4)} in \file{common/parallel.py} runs the calls of a
step a few at a time in threads and returns the results in the order of the items. A step
spends nearly all its time waiting for the model, so four calls at once shorten a report from
tens of minutes to a few. Each thread closes its own database connection when it is done. If a
call fails, the others are still allowed to finish, and the first exception is raised after
all have settled.

\code{shared=} marks the part of the input that is the same in every call of a step: in a
report's gap analysis, the company's profile and its documents. \code{run_prompt} passes it
apart from the data, and the assistant service places it behind the instructions, where the
provider caches it. To profit from that, the first call must have finished before the others
start --- calls started together would each pay for writing the block. \code{shared_first}
in \file{analysis/pipeline.py} does exactly that: it runs the first item alone and hands the
rest to \code{parallel_map}. \Cref{sec:models:cache} gives the measurement behind this.

\subsection{Retries, patience and the errors a caller sees}\label{sec:models:errors}

A failure is handled at three levels, each with a different horizon. The SDK inside the
assistant service repeats a failed provider call twice within seconds. The API's client then
decides between a brief retry and a long wait. What is left reaches the calling code as
\code{GatewayError}.

The client distinguishes two kinds of caller. A call a person waits for --- a citation check
from the corpus page, a test run --- must fail soon enough to tell that person. A call made for
a job step is background work nobody watches: for it a rate limit is \enquote{no failure at
all}, since the provider counts tokens per minute and the limit clears by itself.
\Cref{tab:models:retries} gives the rules of \code{_post}.

\begin{table}[H]
\caption{What the API's client does with a failed call. $n$ counts the retries already made;
\code{generate} is called with one retry by default, \code{embed} with three, the test run of
the prompt editor with none.}\label{tab:models:retries}
\begin{tabularx}{\textwidth}{@{}L{44mm}Y@{}}
\toprule
\hd{Failure} & \hd{What the client does} \\
\midrule
A rate limit (429), patient call & Waits and asks again, without using up a retry: as long as the provider asked (\code{retryAfterS}), otherwise \code{RATE_LIMIT_WAIT_S} = 20~s, then 40, then 60 at most; each wait a quarter more or less at random, so that the calls of one step do not all return in the same second. After \code{RATE_LIMIT_PATIENCE_S} = 900~s of waiting in all, the limit is treated like the next row. \\
A rate limit (429), other calls; 408, 409, 425, 500, 502, 503, 504 & Retried up to the caller's number of retries, after $\min(2\cdot 3^{n},\,30)$ seconds: 2, 6, 18. \\
The connection failed (\code{GATEWAY_UNREACHABLE}) & Retried at least \code{CONNECTION_RETRIES} = 5 times when the caller allows any retry (after 2, 6, 18, 30 and 30~s): the assistant service was restarting, the call itself was sound. \\
No answer in time (\code{GATEWAY_TIMEOUT}) & Retried up to the caller's number of retries. \\
Any other 4xx, a spent budget (402) included & Raised at once. Asking again cannot help. \\
\bottomrule
\end{tabularx}
\end{table}

\Cref{tab:models:errors} lists everything a caller can receive. The assistant service builds
these answers in \code{providerFailure} (\file{routes/internal.js}) from the classification in
\file{middleware/errors.js}: a spent budget becomes a 402, a rate limit a 429 with the wait
the provider named, and anything else keeps the provider's own status, or 502 when there was
none --- so that the client can tell a failure worth repeating from one that is not.

\begin{xltabular}{\textwidth}{@{}L{41mm}L{50mm}Y@{}}
\caption{The errors of a model call as the API receives them}\label{tab:models:errors}\\
\toprule
\hd{Status and code} & \hd{What happened} & \hd{What follows} \\
\midrule
\endfirsthead
\toprule
\hd{Status and code} & \hd{What happened} & \hd{What follows} \\
\midrule
\endhead
\bottomrule
\endlastfoot
402 \code{AI_BUDGET_EXHAUSTED} & The gateway refused the call because the budget of its key is spent (\cref{sec:models:budgets}); also a provider's worded quota refusal. & Asked once by the SDK, never retried by the client. A report fails with a message that begins \code{AI_BUDGET_EXHAUSTED}, which the console shows in words, and keeps the work already paid for. \\
429 \code{GENERATION_FAILED} or \code{EMBEDDING_FAILED} & The provider's rate limit outlasted the SDK's two retries, or arrived inside a stream. \code{retryAfterS} carries the wait when the provider named one. & A job step waits it out for up to fifteen minutes; a call a person waits for is retried briefly (\cref{tab:models:retries}). \\
400 \code{UNKNOWN_MODEL} & The model id names no model, an unknown provider, or a provider without credentials. & Raised at once. The step fails with the message. \\
400, a code that begins \code{INVALID_} & A field of the request failed validation (\cref{tab:models:request}). & Raised at once; a defect of the calling code. \\
4xx \code{GENERATION_FAILED} & The provider refused the request itself: a schema that is not strict, a model the key may not use, an input longer than the model takes. & Raised at once, with the provider's words in the message. \\
5xx \code{GENERATION_FAILED} & The provider's own error, after the SDK's retries. & A 500, 502, 503 or 504 is retried by the client, then raised. \\
502 \code{GENERATION_FAILED} or \code{EMBEDDING_FAILED} & No provider status: the answer was not valid JSON or did not match the schema, the call ran into its time ceiling, or the network to the provider failed. & Retried by the client, then raised. \\
0 \code{GATEWAY_TIMEOUT} & The assistant service did not answer within the client's timeout. & Retried, then raised. \\
0 \code{GATEWAY_UNREACHABLE} & No connection to the assistant service. & Retried five times (not at all for a caller that allows no retry), then raised. \\
403 \code{INTERNAL_ONLY} & The shared secret is missing or wrong. & Raised at once; a defect of the configuration. \\
\code{PromptNotConfigured} & Not an error of the call: the prompt has no active version. & Raised by \code{run_prompt} before any call. \\
\end{xltabular}

What a caller does with a raised error depends on whether it can do without the model. The
steps of a report cannot: the exception ends the job, and the report records the message
(\code{failure_message} in \file{analysis/pipeline.py}; \cref{ch:report}). Several smaller
readers can, and each has a stated fallback:

\begin{itemize}
  \item the citation check judges by word overlap instead of the model
    (\file{corpus/citations.py}) --- except that inside a report a 402 is raised like any other
    failure of a step, so that a report made without its budget is not mistaken for a
    complete one;
  \item the reading of the structure of a standard falls back to the structure rules alone
    and, when they find nothing, to cutting the standard by pages (\code{read_standard} in
    \file{corpus/outline.py});
  \item the reading of an act's scope and of its sanction provisions logs the failure and
    leaves the act unread, to be asked again by a later run (\file{corpus/scope.py},
    \file{corpus/sanctions.py});
  \item a search whose query cannot be embedded runs on keywords alone
    (\file{corpus/search.py}, \file{companies/search.py});
  \item the test run of the prompt editor answers 502~\code{MODEL_ERROR} with the cause.
\end{itemize}

The chapters on each path say what the fallback means for the result.


\section{Usage: what is recorded for every call}\label{sec:models:usage}

Every model call leaves one row in the table \code{llm_calls}. The row is written by
\code{record_call} in \file{assistant/usage.py}, which four kinds of caller use:
\code{run_prompt} for every structured generation, the test run of the prompt editor
(\code{test_run} in \file{prompts/api/views.py}), the embedding callers for every batch, and
the internal route \route{POST}{/internal/assistant/turn-log}, through which the assistant
service logs each chat turn and each conversation title it wrote.

\subsection{Five numbers}

What a call costs depends on more than how many tokens went in, because the provider bills
three kinds of input at three prices (\cref{sec:models:cache}). A call is therefore counted in
five numbers, read from what the SDK reports by \code{usageOf} in \file{llm/usage.js}:

\begin{enumerate}
  \item \code{inputTokens} --- the whole input;
  \item \code{cachedInputTokens} --- the part of the input the provider read from its prompt
    cache;
  \item \code{cacheWriteTokens} --- the part of the input the provider wrote to that cache;
  \item \code{outputTokens} --- the whole output;
  \item \code{reasoningTokens} --- the part of the output that was the model's reasoning.
\end{enumerate}

A provider that reports no cache leaves the second and third at zero. For a chat turn, which
consists of several model steps, the five numbers are the sums over the steps
(\code{addUsage}).

\subsection{The cost of a call}

\code{cost_of} turns the counts into US dollars. With $I$ the input tokens, $c$ of them read
from the cache, $w$ written to it and $O$ the output tokens, and with the model's four prices
per token --- $p$ for input, $p_r$ for input read from the cache, $p_w$ for input written to
it, $p_o$ for output ---
\[
  \mathrm{cost} \;=\; p\,(I - c - w) \;+\; p_r\,c \;+\; p_w\,w \;+\; p_o\,O ,
\]
rounded to a millionth of a dollar. Three details keep the figure honest. The two parts are
clamped so that they can never add up to more than the input ($c \le I$, $w \le I - c$). A
model for which the price list names no cache price is charged the plain input price for that
part. And a model with no price at all costs 0, which is logged once per model --- the usage
row is still written, with its tokens.

An example with the prices of \model{gpt-6.1-sol} (a test in \file{tests/test_prompts.py}
pins the same arithmetic): a call with 60\,000 input tokens, of which 55\,000 were
read from the cache and 4\,000 written to it, and 3\,000 output tokens costs
$1\,000 \times 2 + 55\,000 \times 0.10 + 4\,000 \times 2.50 + 3\,000 \times 10$ millionths of
a dollar, that is \$0.0475. Priced by its token counts alone it would have been recorded as
\$0.15.

\subsection{The price list}

Prices are not written into the code as the truth; they are read from the gateway.
\code{price_table} asks the assistant service for the model list, which carries for each model
the price per input and per output token and, where the gateway names them, the two cache
prices. The table is kept in memory for an hour (\code{PRICE_TTL_S} = 3600; the assistant
service itself caches the gateway's list for ten minutes). Three layers stand behind it, so
that a hiccup of the gateway neither loses a usage row nor zeroes a budget:

\begin{enumerate}
  \item every successful fetch is also stored in the table \code{model_prices} (one row per
    model: the four prices and the time);
  \item a failed fetch falls back to that stored table, which then serves for the hour like a
    fresh one; only when nothing is stored either is the fetch tried again after five minutes
    (\code{PRICE_RETRY_S} = 300);
  \item under both lies \code{KNOWN_PRICES}, the prices known when the code was written
    (\cref{tab:models:prices}), for an installation that has never reached the gateway and
    for a model the fetched list does not name.
\end{enumerate}

\begin{table}[H]
\caption{The prices the code ships as its last resort (\code{KNOWN_PRICES} in
\file{assistant/usage.py}), in US dollars per million tokens. They are the gateway's prices of
27 September 2026; those of the GPT-6 models and all cache prices are of 10 November 2026. A
dash means the list names no such price and the input price applies.}\label{tab:models:prices}
\begin{tabularx}{\textwidth}{@{}Yrrrr@{}}
\toprule
\hd{Model} & \hd{Input} & \hd{Output} & \hd{Read from cache} & \hd{Written to cache} \\
\midrule
\model{azure/gpt-6-astra} & 10 & 50 & 1.00 & 12.50 \\
\model{azure/gpt-6.1-sol} & 2 & 10 & 0.10 & 2.50 \\
\model{azure/gpt-6-sol} & 2 & 10 & 0.20 & 2.50 \\
\model{azure/gpt-5.5} & 5 & 30 & 0.50 & --- \\
\model{azure/gpt-5.6-sol} & 4 & 20 & --- & --- \\
\model{azure/gpt-5.6-terra} & 2 & 12 & 0.20 & 2.50 \\
\model{azure/gpt-5.6-luna} & 0.20 & 1.20 & 0.02 & 0.25 \\
\model{azure/text-embedding-3-large} & 0.13 & 0 & --- & --- \\
\model{azure/text-embedding-3-small} & 0.02 & 0 & --- & --- \\
\bottomrule
\end{tabularx}
\end{table}

The estimate shown before a report starts (\file{analysis/estimate.py}; \cref{ch:report}) uses
the same \code{cost_of} and the same table, for the model of each step's active prompt
version. It prices the input of a report's steps as written to the cache --- what the provider
charges for input it sees for the first time --- and the gap analysis as one call that writes
the shared block and the rest that read it; the one-time extraction of an act never read is
priced at the plain input price. When a model of the run has no price, the estimate
says so (\code{pricesKnown} is false) instead of presenting a total that is only a floor.

\subsection{The row}

\Cref{tab:models:llmcalls} lists every field of a row. Two of them need a word.

\code{kind} is one of four values, enforced by a constraint of the database: \code{chat},
\code{job_step}, \code{embed}, \code{prompt_test}. \code{run_prompt} always writes
\code{job_step}, also for the few calls it makes outside a job (a citation check from the
corpus page); such a row simply has no job step.

\code{purpose} says what the call was for, in at most 80 characters. For a structured
generation it is the slug of the prompt, unless the caller names another (the citation judge
writes \code{citations.check}; the evaluation suites write \texttt{eval:} and the suite's name,
except the citations suite, whose calls go through the citation judge and carry its purpose).
For embeddings it is one of the four purposes of \cref{tab:models:embedcallers}; for the chat
it is \code{chat} for a turn and \code{chat_title} for a title.

\begin{xltabular}{\textwidth}{@{}L{38mm}Y@{}}
\caption{The fields of a row of \code{llm_calls} (\file{assistant/models.py}, \code{LlmCall})}\label{tab:models:llmcalls}\\
\toprule
\hd{Field} & \hd{Holds} \\
\midrule
\endfirsthead
\toprule
\hd{Field} & \hd{Holds} \\
\midrule
\endhead
\bottomrule
\endlastfoot
\code{id} & The row's own identifier. \\
\code{kind} & \code{chat}, \code{job_step}, \code{embed} or \code{prompt_test}. \\
\code{purpose} & What the call was for: a prompt's slug or one of the fixed purposes above. \\
\code{user} & The user the call is charged to; empty for shared work (\cref{tab:models:charged}). The daily budget sums over this field. \\
\code{thread_id} & For a chat turn, the conversation. \\
\code{job_step_id} & For background work, the step of the job. When a step ends, the worker adds up the step's rows --- calls, tokens, cost, by model --- into the step's own record (\file{jobs/runner.py}), which is where a report's cost per step comes from. \\
\code{prompt_version_id} & The prompt version whose instructions ran; empty for embeddings and for a test run of an unsaved draft. \\
\code{model} & The model that ran, in the \texttt{provider:model} form; for a call through the generate or embed route, as the assistant service named it in its answer. \\
\code{input_tokens}, \code{cached_input_tokens}, \code{cache_write_tokens}, \code{output_tokens}, \code{reasoning_tokens} & The five numbers. \\
\code{cost_usd} & The cost computed by \code{cost_of} when the row was written, to six decimals. A later change of prices does not alter it. \\
\code{duration_ms} & How long the call took. \\
\code{tool_calls} & For a chat turn, the list of tools the model called. \\
\code{outcome} & \code{ok}, or for a chat turn how it ended otherwise (for example \code{timeout}, \code{budget}, \code{error}); for a title, \code{empty} when the model wrote none. \\
\code{transcript} & For a chat turn: the user's message (at most 20\,000 characters) and the assistant's answer (at most 100\,000); for a title, the opening of the first message and the title written. Empty for every other call. \\
\code{created_at} & When the row was written (UTC). \\
\end{xltabular}

\paragraph{The transcript.} The chat is the one feature whose calls keep their content in the
usage row: for a turn, what the user asked and what the assistant answered; for a title, what
the model read and wrote. The header of the model calls
it \enquote{the audit of what the assistant told whom}. It is served only in the detail of a
call (\route{GET}{/api/usage/calls/{id}}), which a user may read for their own calls and an
administrator for any. The inputs and outputs of pipeline calls are not copied into
\code{llm_calls}; their results live in the tables of each path.

\paragraph{Telemetry never fails the caller.} \code{record_call} runs its insert in a
savepoint of its own and catches every exception: a failure to write the row is logged and the
model's answer is still used.

\paragraph{What is not recorded.} Two things leave no row. The first is deliberate and stated
in \file{docs/markdown/en/architecture.md}: the embedding of a search query, a few tokens per search, is
made without a usage row. The second follows from the order of the code: \code{run_prompt} and
the embedding callers record a call after its answer has arrived, so a call from the worker
that fails --- including one whose answer the schema check rejected --- has no row, although the
provider may have billed its tokens. A chat turn that fails is recorded, with its
\code{outcome}.

The usage page of the console reads these rows through \route{GET}{/api/usage} and
\route{GET}{/api/usage/calls}: cost per day, by kind, by model and --- for an administrator ---
by user (\cref{ch:console,ch:api}).


\section{Budgets}\label{sec:models:budgets}

Two budgets limit what the system can spend. One belongs to the gateway and covers the whole
installation; the other is the system's own and covers one user for one day.

\subsection{The team budget of the gateway key}

The key the assistant service calls the gateway with carries a \term{team budget}, set and
counted at the gateway. Nothing in this system can change it. When it is spent, the gateway
answers every call with HTTP~402, and the system says so plainly instead of failing
generically: the assistant service turns the refusal into 402~\code{AI_BUDGET_EXHAUSTED} on
the internal door and into the same code in the chat; a report that fails this way records a
message that begins with that code and keeps everything it had already paid for --- the
inventories of documents and the extracted obligations are not bought again by the next run.

Two details of the classification (\code{classifyProviderError} in
\file{middleware/errors.js}) are worth knowing. A 402 is always a spent budget; a 400, a 403, a
429 or an error without a status counts as one only when its text contains both a word for money (budget, quota, billing,
credit balance) and a word for running out. And rate-limit wording never reads as a spent
budget, even when it mentions a quota: the two demand opposite reactions --- wait, or stop.
\code{finalBudgetRefusals} then marks a budget refusal as not retryable, because the SDK
would otherwise repeat a 429 that no wait of seconds can cure.

\file{docs/markdown/en/operations.md} records that the team budget was \$10 on 27 September 2026, the day
it was reached.

\subsection{The daily budget of a user}

The system's own budget \enquote{only keeps one account from draining the shared key}, as
\file{docs/markdown/en/operations.md} puts it. Each user may spend a fixed amount per UTC day:

\begin{itemize}
  \item the default is the setting \code{LLM_DAILY_BUDGET_USD}, \$10.00 since 20~November
    2026. It was \$2.00, set in September, when a chat turn cost cents. On the conversation
    model in use since November the turn that takes in a company's documents costs about
    \$2.90 and a later turn of a conversation that has read them about \$1.00 (the median of
    the live system's turns of 17~November was \$1.02), so the first hand-over of documents
    spent an ordinary account's day and its next question was refused;
  \item an administrator may give an account its own amount (\code{daily_budget_usd} on the
    user; in the console, on the \ui{Naudotojai} (Users) page): empty means the default, 0
    means unlimited, otherwise an amount up to \$10\,000;
  \item an administrator's own account is unlimited unless an amount is set for it.
\end{itemize}

\code{budget_status(user_id)} in \file{assistant/usage.py} computes the state: the limit, the
amount used --- the sum of \code{cost_usd} over the user's rows in \code{llm_calls} since
00:00~UTC --- and what remains (nothing, when the limit is 0, which means unlimited). Because
it sums over usage rows, the budget counts exactly what \cref{tab:models:charged} says is
charged to a user, and nothing else.

The budget is checked \emph{before} work starts, never in the middle of it
(\cref{tab:models:budgetchecks}). A report or a work product must have \$0.25 left to start
(\code{MIN_BUDGET_TO_START_USD}); once started, a run is never stopped halfway, even if it passes the limit. The estimate a user confirms before the
start (\cref{ch:report}) is what tells them the size of the run.

\begin{table}[H]
\caption{Where the daily budget is checked, and what happens when it is spent}\label{tab:models:budgetchecks}
\begin{tabularx}{\textwidth}{@{}L{31mm}L{50mm}Y@{}}
\toprule
\hd{Feature} & \hd{Checked in} & \hd{When the budget is spent} \\
\midrule
A chat turn & assistant service, \file{routes/chat.js}, \code{checkBudget}, before every turn (it asks the API on \code{/internal/assistant/budget}) & The turn is refused with 429~\code{BUDGET_EXCEEDED} when nothing remains. If the check itself fails, the turn goes through. \\
Starting a report & \file{analysis/api/views.py}, \route{POST}{/api/reports} & Refused with 429~\code{BUDGET_EXCEEDED} when less than \$0.25 remains. \\
Approving the scope of a report & \file{analysis/api/views.py}, \route{POST}{/api/reports/{id}/scope} & The same: the costly second part needs \$0.25 left. \\
Starting a work product & \file{programme/api/views.py}, \route{POST}{/api/companies/{id}/programme} & The same. \\
A test run in the prompt editor & \file{prompts/api/views.py}, \code{test_run} & Refused with 429~\code{BUDGET_EXCEEDED} when nothing remains. \\
A citation check from the corpus page or the chat & \file{corpus/api/views.py}, \route{POST}{/api/corpus/check-citations} & Not refused: the check runs, judged by word overlap instead of the model. \\
Indexing an uploaded document & \file{companies/indexing.py}, the job \code{documents.index} & Nothing is embedded. The job queues itself again for 00:05~UTC of the next day and reports that it is waiting for the budget. \\
\bottomrule
\end{tabularx}
\end{table}

The last row has a reason of its own, given in the code: without the wait, uploading and
deleting a document in a loop would embed, and charge the gateway, without end.

\subsection{Who is charged for what}

Whether a call counts against a user is decided by one argument: the \code{user_id} its caller
passes to \code{run_prompt} or \code{record_call}. The rule behind the choices is that work
done \emph{for one user's company} is that user's, and work done \emph{on the shared legal
corpus} --- which every later report of every company reuses --- is nobody's
(\cref{tab:models:charged}). Shared work still spends the team budget of the gateway key; it
is only not held against a person's day.

\begin{table}[H]
\caption{Who is charged for a model call}\label{tab:models:charged}
\begin{tabularx}{\textwidth}{@{}L{62mm}Y@{}}
\toprule
\hd{Work} & \hd{Charged to} \\
\midrule
A chat turn, and the title written for a new conversation & The chatting user. \\
A report: document inventories, the citation audit, applicability, gap analysis, action plan and summary; the indexing of documents the report finds unindexed & The user who started the report; after a scope review, the calls of the second part are charged to whoever approved the scope (the job of that part is theirs). \\
A company's own requirements, read from its policies and assessed in a report & The user who started the report. \\
A work product & The user who started it, except the citation judge's call of its references step, which is recorded under the job step with no user. \\
Indexing an uploaded document (embeddings) & The uploader. \\
A citation check from the corpus page or the chat & The user who asked. \\
A test run in the prompt editor & The administrator who ran it. \\
Extraction of obligations that a report triggers because an act was never read & Nobody: shared corpus work. \\
Extraction of obligations, or the judge over an act's obligations, queued by an administrator & That administrator. \\
Reading an act's scope and its sanction provisions; reading the structure of a standard & Nobody. \\
Ingesting an act (embeddings of its chunks); the embeddings of obligations & Nobody. \\
The embedding of a search query & Not recorded at all. \\
\bottomrule
\end{tabularx}
\end{table}

\begin{humanbox}{What an administrator decides}
No budget is set or lifted by the system itself. An administrator decides the daily amount
of each account and may lift it; the operators of the gateway decide the team budget. An
administrator also decides, per prompt, which model runs (\cref{sec:models:defaults}) --- the
largest single lever on what a report costs. A user decides to start a run after seeing its
estimated cost.
\end{humanbox}


\section{The prompt cache}\label{sec:models:cache}

The provider behind the gateway keeps a \term{prompt cache}: when a request opens with text
it has already processed, it reads that opening from the cache and bills it at a
fraction of the input price. The system did not choose this, but its design had to follow it,
because the largest calls of a report repeat most of their input.

\subsection{What is cached and what is not}

What exactly the provider caches was established by measurement, and the result is narrower
than \enquote{whatever repeats}.

\begin{lessonbox}{Measured, 10--11 November 2026, on gpt-6.1-sol}
The same block of 8\,800 tokens was sent twice in a row in several positions. It was
\textbf{read from the cache} the second time when it stood in the instructions of the
request. It was \textbf{written anew} --- billed in full --- when it stood at the opening of a
user message whose end differed; when it was a message of its own before a different one;
when a cache key was sent with it; and when a retention hint was sent with it. A request that
repeated or extended an earlier request word for word from its start read what it repeated
from the cache.

So the provider caches two things: the \emph{instructions} of a request, and a \emph{whole
earlier request} that a later one repeats or extends. It caches nothing of a message that
only begins like an earlier one.

The prices on that day for \model{gpt-6.1-sol}: input read from the cache \$0.10 per million
tokens, a twentieth of the input price; every other input token billed as written to the
cache at \$2.50, a quarter above the input price of \$2.
\end{lessonbox}

\subsection{Three consequences for the design}

\paragraph{1. What the calls of a step share goes behind the instructions.} The gap analysis
of a report sends the company's documents with every call. Had they stayed in the user
message, in front of the obligations that differ from call to call, each call would have paid
for them in full. They are therefore sent apart, as \code{shared}, and the generate route
appends them to the instructions --- the one place the provider caches
(\cref{fig:models:anatomy}).

\paragraph{2. The first call of such a step runs alone.} The cache holds the block only once
a call that carried it has been processed. \code{shared_first} runs one call, waits for it,
and only then starts the others four at a time. The applicability step and the gap analysis
of a report use it, and so does the assessment of a company's own requirements.

\paragraph{3. Usage is recorded and priced part by part.} This is the reason for the five
numbers of \cref{sec:models:usage}. Until 11 November 2026 every input token was priced at
the plain input price; \file{docs/markdown/en/architecture.md} records that a chat turn, which sends its
instructions again at each of its steps, was then recorded at up to twice its real cost. The
budgets and the report's cost per step rest on the corrected figures.

\begin{figure}[H]
\centering
\begin{tikzpicture}[x=1mm, y=1mm]
  \tikzset{seg/.style={draw=ink, line width=.6pt, minimum height=9mm, align=center,
                       inner sep=2pt, font=\sffamily\footnotesize, text=body}}
  % the parts of the request
  \node[seg, fill=paper, minimum width=24mm, anchor=west] (i0) at (26,0) {instructions};
  \node[seg, fill=paper, minimum width=66mm, anchor=west] (s0) at (50,0) {shared block: the profile and the documents};
  \node[seg, fill=paper, minimum width=30mm, anchor=west] (d0) at (116,0) {data of this call};
  \node[note, anchor=east] at (24,0) {one request};
  \draw[decorate, decoration={brace, amplitude=4pt}, draw=muted] (26,6.5) -- node[note, above=4pt, align=center] {sent as the instructions --- the part the provider caches} (116,6.5);
  \draw[decorate, decoration={brace, amplitude=4pt}, draw=muted] (116,6.5) -- node[note, above=4pt, align=center] {the user message} (146,6.5);
  % the first call
  \node[seg, fill=warnbg, minimum width=120mm, anchor=west] (a) at (26,-17) {written to the cache: about 59\,000 tokens at \$2.50 per million};
  \node[note, anchor=east, align=right] at (24,-17) {the first call,\\alone};
  % later calls
  \node[seg, fill=goodbg, minimum width=90mm, anchor=west] (b1) at (26,-31) {read from the cache: about 55\,000 tokens at \$0.10 per million};
  \node[seg, fill=warnbg, minimum width=30mm, anchor=west] (b2) at (116,-31) {written};
  \node[note, anchor=east, align=right] at (24,-31) {each later\\call};
\end{tikzpicture}
\caption{The anatomy of a request that carries a shared block, with the figures of the
sample company's gap analysis (26 policies, 261\,000 characters; measured on 10--11 November
2026). The instructions and the shared block are identical in every call of the step and are
read from the provider's cache once the first call has written them; the data of each call is
new input every time.}\label{fig:models:anatomy}
\end{figure}

\subsection{The measured effect}

On the sample company every gap-analysis call carries about 59\,000 input tokens, nearly all
of them the company's 26 policies. \Cref{fig:models:cachecost} shows what one call cost in
the three situations, as \file{docs/markdown/en/architecture.md} records them for 10--11 November 2026 on
\model{gpt-6.1-sol}.

\begin{figure}[H]
\centering
\begin{tikzpicture}
\begin{axis}[
  xbar, width=92mm, height=50mm,
  xmin=0, xmax=0.22, bar width=6.5mm,
  enlarge y limits=0.28,
  axis x line*=bottom, axis y line*=left,
  axis line style={draw=muted}, tick style={draw=muted},
  ytick=data, symbolic y coords={before,read,write},
  yticklabels={{documents in the user message\\(before the change)},
               {a later call: the block\\read from the cache},
               {the first call: the block\\written to the cache}},
  yticklabel style={align=right, text width=50mm, font=\sffamily\footnotesize, text=body},
  xtick={0,0.05,0.10,0.15,0.20},
  xticklabel style={font=\sffamily\footnotesize, text=muted, /pgf/number format/fixed,
                    /pgf/number format/precision=2, /pgf/number format/fixed zerofill},
  scaled x ticks=false,
  xlabel={US dollars for one call}, xlabel style={font=\sffamily\footnotesize, text=muted},
  nodes near coords={\$\pgfmathprintnumber[fixed, precision=3]{\pgfplotspointmeta}},
  nodes near coords style={font=\sffamily\footnotesize, text=body, anchor=west},
  point meta=x]
  \addplot[fill=mid, draw=ink] coordinates {(0.17,before) (0.045,read) (0.18,write)};
\end{axis}
\end{tikzpicture}
\caption{The cost of one gap-analysis call of the sample company (about 59\,000 input
tokens, \model{gpt-6.1-sol}, measured on 10--11 November 2026). The first call of the step
pays \$0.18 to write the shared block; each later call reads 55\,000 tokens of it from the
cache and costs about \$0.045. Before the change, with the documents in the user message,
every call cost \$0.17.}\label{fig:models:cachecost}
\end{figure}

\file{docs/markdown/en/architecture.md} extrapolates, from a run of 11 November 2026 on a copy of the
database with every act's obligations cut to its first fifteen, that the full report then
came to about 82 such calls and roughly \$5, against \$16.04 for the report made the day
before; other changes to the pipeline between the two runs share in that difference
(\cref{ch:report}). The first full report on the live system, on 13 November 2026 over a
corpus grown to 68 enabled acts, made 154 gap-analysis calls: they sent 9.12 million input
tokens, of which 8.44 million (92.5\,\%) were read from the provider's cache and 0.68 million
written to it. The same mechanism serves the chat without any arrangement:
each step of a turn extends the request of the step before, so a step reads its instructions
and the earlier steps from the cache and pays in full only for what the last step added
(\cref{ch:chat}).

The cache is the provider's and nothing in the system depends on it for correctness. A call
that finds the cache empty returns the same answer and costs more; the usage row shows which
of the two happened.


\section{Rate limits}\label{sec:models:ratelimits}

Apart from money, the provider limits throughput: it counts the tokens a key sends per
minute. A report's many large calls can reach that limit within minutes.

\begin{lessonbox}{A failure, 10 November 2026}
Four gap-analysis calls of 59\,000 tokens each, running at a time, met the provider's
tokens-per-minute limit within two minutes. The client's brief retries --- after 2, 6 and
18~seconds --- met the limit again, and the report failed. Since then a call made for a job
step waits a rate limit out, as long as the provider asks and for up to fifteen minutes in all
(\code{RATE_LIMIT_PATIENCE_S}), instead of failing its step.
\end{lessonbox}

\paragraph{How a rate limit arrives.} It comes in one of two shapes, and
\code{providerPause} in \file{middleware/errors.js} reads both. The ordinary shape is an HTTP
429 that is not a budget refusal, with or without a \code{Retry-After} header. The other was the
one met on 10 November: the provider had already begun its stream and reported the limit as
an error event inside it, with no HTTP status at all; the SDK hands that on as a failure to
parse the stream. Such an error is recognised by its words (\enquote{rate limit},
\enquote{call rate}, \enquote{retry after} and a number). The wait is taken from the header
when there is one, else from the seconds the message names, else it is unknown.

\paragraph{What is done with it.} The SDK inside the assistant service first repeats the call
twice. A limit that outlasts that, or one that arrived inside a stream, is answered to the API
as a 429, with the field \code{retryAfterS} when the provider named a wait. From there the rules of \cref{tab:models:retries}
apply: a job step waits and asks again without counting it as a failure; a call a person
waits for is retried briefly and then fails with a message. In the chat a rate limit is shown
as a notice to try again (\cref{ch:chat}).

A step of a report runs four calls at a time (\code{WORKERS} in
\file{analysis/pipeline.py}); it is the patience of each call that absorbs the limit.

\paragraph{The system's own limit on chat turns.} One rate limit is the system's own and
concerns people, not tokens: a user may start at most 12 chat turns a minute
(\code{RATE_LIMIT_TURNS}, a fixed window of \code{RATE_LIMIT_WINDOW_MS} = 60\,000~ms kept in
memory by \file{services/ratelimit.js}). A turn over the limit is refused with
429~\code{RATE_LIMITED} and the seconds left in the window (\cref{ch:chat}).


\section{The default model of every step}\label{sec:models:defaults}

Which model a prompt runs on is a field of its active version, and an administrator can
change it in the prompt editor (\cref{ch:prompts}). The versions the system ships with name
three language models and one embedding model. \Cref{tab:models:defaults} gives the default of
every prompt and the recorded reason for it. Each reason is a dated observation on the sample
company's material; none is a general ranking of models.

\begin{xltabular}{\textwidth}{@{}L{36mm}L{40mm}Y@{}}
\caption{The default model of every prompt, with the recorded reason. Prices are US dollars per million input and output tokens.}\label{tab:models:defaults}\\
\toprule
\hd{Prompts} & \hd{Default model} & \hd{Why} \\
\midrule
\endfirsthead
\toprule
\hd{Prompts} & \hd{Default model} & \hd{Why} \\
\midrule
\endhead
\bottomrule
\endlastfoot
\prompt{chat.assistant} & \model{gpt-6-astra} (10 / 50) & The shipped default of the chat; the model before it was \model{gpt-5.5}. The repository's documents record no comparison for the chat. A user can pick another model for a conversation (\cref{ch:chat}). \\
\prompt{chat.title} & \model{gpt-5.6-luna} (0.20 / 1.20) & A few words written from the opening of a conversation's first message: the small model. \\
The report: \prompt{report.inventory}, \prompt{report.applicability}, \prompt{report.gap_analysis}, \prompt{report.action_plan}, \prompt{report.plan_merge}, \prompt{report.summary}; a company's own requirements: \prompt{company.internal_requirements}, \prompt{report.internal_gap} & \model{gpt-6.1-sol} (2 / 10) & Measured on 10 November 2026 on the sample company's real requests --- three gap-analysis batches of fifteen obligations and two applicability batches of four acts: \model{gpt-6.1-sol} and \model{gpt-6-astra} gave the same status for 43 of 45 obligations and the same verdict for all eight acts, at a fifth of the price; \model{gpt-5.6-terra} differed on five obligations, each time by overlooking a passage or by holding a policy to something the obligation does not ask. An action plan and a summary written by each from the same twenty open findings covered all twenty. \\
\prompt{corpus.extract_obligations}, \prompt{corpus.extract_requirements} & \model{gpt-6.1-sol} & One-time corpus work that every report reuses, so the most thorough reader measured. On 10 November 2026, on two batches (MiCA Art.~68--75 and Art.~10--13 of PPTFPĮ, the Lithuanian law on the prevention of money laundering and terrorist financing): \model{gpt-5.6-terra} listed 44 and 25 obligations, \model{gpt-6.1-sol} 56 and 32, \model{gpt-6-astra} 55 and 31 at about five times the price. \\
\prompt{corpus.standard_outline}, \prompt{corpus.scope}, \prompt{corpus.sanctions} & \model{gpt-6.1-sol} & One-time corpus work as well; \file{docs/markdown/en/architecture.md} groups the structure of a standard and the scope of an act with extraction. No comparison of their own is recorded for these three. \\
\prompt{corpus.check_citations} & \model{gpt-6.1-sol} & Short comparisons, a few calls a report. The check was first measured on \model{gpt-5.6-luna}: it judged all 12 of the wrong MiCA references of the sample company correctly, where similarity of embeddings did no better than chance. \\
\prompt{corpus.verify_obligations} (the judge) & \model{gpt-6.1-sol} & One-time corpus work: judging the whole register of 5\,960 obligations took 369 calls and \$9.80 on 13 November 2026. \\
The six prompts of the work products, whose slugs begin \prompt{programme.} & \model{gpt-6.1-sol} & Two calls a product over the company's whole library; the second reads the library from the cache. A product costs about \$0.4--0.5 (\file{docs/markdown/en/architecture.md}). \\
Embeddings (no prompt) & \model{text-embedding-3-large}, 1536 dimensions (0.13) & The comparison of 28 September 2026 (\cref{tab:models:embedcompare}). \\
\end{xltabular}

\paragraph{One model for the whole report.} \file{docs/markdown/en/architecture.md} draws the conclusion
from the measurement of 10 November: \model{gpt-6-astra} costs five times
\model{gpt-6.1-sol} for answers that agreed --- about \$57 against \$11 for the gap analysis
of the sample company's 2\,200 obligations --- so every step of a report runs on
\model{gpt-6.1-sol}. The same document notes the price of that choice in time: the GPT-6
models wrote about 70 tokens a second on that day, half the speed of \model{gpt-5.6-terra},
which is why a dense batch takes four to six minutes and why the ceilings of
\cref{tab:models:timeouts} are what they are.

\paragraph{The embedding model.} Two models of the gateway were compared on 28 September 2026
on the system's own data (\cref{tab:models:embedcompare}). The larger one is clearly better
exactly where the system needs it: at finding a Lithuanian provision from an English text.
The chat often searches in English and companies' documents are often English, while the
corpus is Lithuanian. \file{docs/markdown/en/architecture.md} puts the cost of the larger model at about
\$0.13 per million tokens, some \$0.30 for embedding everything again.

\begin{table}[H]
\caption{The two embedding models compared, 28 September 2026 (from \file{docs/markdown/en/architecture.md}). The first three rows count how often an obligation's text found its own article among the Lithuanian chunks of the act, out of 40.}\label{tab:models:embedcompare}
\begin{tabularx}{\textwidth}{@{}Yrr@{}}
\toprule
\hd{Test} & \hd{-large} & \hd{-small} \\
\midrule
English obligation text, MiCA & 38 & 23 \\
English obligation text, DORA & 33 & 29 \\
English obligation text, PPTFPĮ & 37 & 34 \\
Lithuanian queries (MiCA / DORA / PPTFPĮ) & 39 / 35 / 37 & 36 / 37 / 36 \\
Company evidence: the right document among an obligation's six nearest passages, of 41 & 38 & 35 \\
The same, mean reciprocal rank & 0.83 & 0.79 \\
\bottomrule
\end{tabularx}
\end{table}

What the models are told on each of these prompts, and how their answers are checked beyond
their shape, is the subject of the chapters that follow; \cref{app:prompts} lists every
prompt, and \cref{ch:control} gathers the costs measured for a whole report.
